\section{Artifact Interface and Validation}

\tool treats a tokenizer export as an inspectable resource artifact, not as an
opaque component of a trained recommender. The minimum input is a table
\texttt{sid\_assignments(item\_id, sid\_0, ..., sid\_L)}. Optional tables provide
metadata, interactions, refresh pairs, and generator outputs. This
mapping-first contract lets adapters normalize heterogeneous
repositories while keeping the inspection independent of a particular training loop.
Figure~\ref{fig:audit-sid-pipeline} summarizes the architecture and artifact
boundary.

\begin{figure}[t]
\centering
\includegraphics[width=\columnwidth]{figures/fig1_audit_sid_pipeline.pdf}
\caption{\tool architecture.}
\Description{SIDInspector architecture. A SID export passes through adapter,
normalized contract, validator, and probe engine. Optional metadata,
interactions, refresh mappings, and traces attach to the contract. Outputs are
D reports, table snapshots, a reproducibility matrix, and extension hooks.}
\label{fig:audit-sid-pipeline}
\end{figure}

\paragraph{Artifact contract.}
The required contract has three checks: each row has a stable item key, every
\sid level is a discrete token, and the diagnostic item universe is joinable to
metadata and interaction slices. Metadata enables semantic/category slices;
interactions enable collaborative-neighborhood and head-tail diagnostics. The
\texttt{generator\_outputs} table is optional because many public tokenizer
releases expose item codes but no trained generator traces.

\paragraph{Adapter specification.}
An adapter is a thin export layer, not a retraining recipe. It must emit
\texttt{sid\_assignments} with stable item keys, a method/dataset label, and
one column per \sid level. If available, it also emits metadata, interactions,
refresh-pair snapshots, or generator traces using the same item keys. The
validator checks key uniqueness, missing mappings, depth consistency, join
coverage, and source provenance before any row can enter the diagnostic tables.
The minimal schema is deliberately small:
\begin{center}
\small
\texttt{item\_id}, \texttt{method}, \texttt{dataset}, \texttt{sid\_0}, \ldots,
\texttt{sid\_L} $\rightarrow$ \{\texttt{D1}, \ldots, \texttt{D5}\}.
\end{center}
Interactions add D3/D4 slices; paired mappings add D6; generator traces add
D7. This makes a new tokenizer easy to attach without rewriting D1--D5.
Current D1--D5 reports use a rectangular code table; variable-length tokenizers
can be attached by exporting padded or masked levels plus realized length, but
variable-length serving behavior remains structural D5 evidence unless
generator traces are available.

\paragraph{Terminology.}
A \emph{diagnostic probe} is a named metric family: D1 utilization, D2
aliasing, D3 neighborhood alignment, D4 popularity allocation, D5 structural
cost, D6 temporal churn, or D7 generation traces. A \emph{controlled mechanism
probe} is an input row designed to activate one known mechanism and check
whether the diagnostic probe reacts. The main paper tables use D1--D5. D6 is
implemented and released as an extension example for paired refresh mappings;
D7 is only an interface hook until generator outputs or beam
traces are supplied.

\paragraph{Claim scope.}
\tool reports tokenizer artifact rows, reference rows, and mechanism probes
under the same schema. The current tokenizer lines are GRID-style RQ-KMeans,
ReSID/GAOQ, LETTER, and LC-Rec. The first two support the same-item diagnostic
case study; the latter two test released item-index adapters. References and
probes calibrate the diagnostics rather than increasing tokenizer coverage. Optional
rows exercise extensions such as drift or non-Amazon schemas, while
literature-only methods motivate future adapters.

\paragraph{Validation before metrics.}
Adapter outputs are interpreted only after checks for item-count agreement,
missing codes, duplicate keys, depth consistency, and provenance. This guard is
necessary because public SID repositories may expose checkpoints, item mappings,
codebooks, intermediate feature files, or only partial releases. Once the
normalized tables pass validation, D1--D5 can be compared across artifacts.

Prior recommender resources such as RecList and Elliot motivate diagnostic and
reproducible evaluation beyond a single aggregate metric~\cite{chia2022reclist,anelli2021elliot}.
\tool applies that resource logic to a newer artifact boundary: the item
address space consumed by generative recommenders.
